% !TEX program = lualatex
\documentclass[professionalfonts]{beamer}
\newcommand{\slidemode}{1}  % カラー投影モード
\usepackage{beamer_template}
\usepackage{enumerate}

\newcommand{\LectureNum}{13}
\newcommand{\LectureTitle}{フーリエ変換の性質と公式}
\newcommand{\CourseName}{応用数学}
\renewcommand{\TermName}{前期}

\begin{document}

\begin{frame}{第13回　フーリエ変換の性質と公式}
  \begin{block}{本回の内容}
  フーリエ変換の性質と公式
  \end{block}
  \vfill\hfill{\scriptsize \textcolor{gray}{第3章 フーリエ解析　§2}}
\end{frame}

\begin{frame}{フーリエ変換の性質}
  \begin{block}{}
  \begin{description}
    \item[I. 線形性]\mbox{}\\$F[c_{1}f_{1}(x) + c_{2}f_{2}(x)] = c_{1}F_{1}(u) + c_{2}F_{2}(u) \ (c_{1}, c_{2} は定数)$
    \item[II. 移動則（時間シフト）]\mbox{}\\$F[f(x-a)] = e^{-iau} F(u) \ (a は実数)$
    \item[III. 移動則（周波数シフト）]\mbox{}\\$F[e^{iax}f(x)] = F(u-a) \ (a は実数)$
    \item[IV. 相似性（スケーリング）]\mbox{}\\$F[f(ax)] = (1/|a|) F(u/a) \ (a は 0 でない実数)$
    \item[V. 微分則]\mbox{}\\$F[f^{(n)}(x)] = (iu)^n F(u) \ (n は自然数)$
    \item[VI. 双対微分則]\mbox{}\\$F[(-ix)^n f(x)] = F^{(n)}(u) \ (n は自然数)$
  \end{description}
  \end{block}
\end{frame}

\begin{frame}{パーセバルの等式}
  \begin{block}{}
  \[
    \int_{-\infty}^\infty |f(t)|^{2} dt = (1/2\pi) \int_{-\infty}^\infty |F(\omega)|^{2} d\omega
  \]
  \end{block}
\end{frame}

\begin{frame}{例題4　（p.\,99）}
  \begin{exampleblock}{問題}
    次の等式を証明せよ\\[2pt]
    $(1) F[e^{-a^{2}x^{2}} * e^{-b^{2}x^{2}}] = (\pi/(ab)) e^{-u^{2}(a^{2}+b^{2})/(4a^{2}b^{2})}$ 等（たたみこみのフーリエ変換と 98 ページの公式 (1) を利用）\\[2pt]
    (2) 上記の結果と比較して、畳み込み関数の閉じた形を求める\\[2pt]
  \end{exampleblock}
\end{frame}

\begin{frame}{例題4　解答}
  $(1)$ p.\,98 公式 $: F[e^{-cx^{2}}] = \sqrt{\pi/c}\cdot e^{-u^{2}/(4c)}$ より\\[4pt]
  $\hphantom{(1)\ }F[e^{-a^{2}x^{2}}] = (\sqrt{\pi}/a)e^{-u^{2}/(4a^{2})},\ F[e^{-b^{2}x^{2}}] = (\sqrt{\pi}/b)e^{-u^{2}/(4b^{2})}$\\[4pt]
  たたみこみ定理 $: F[f*g] = F(u)G(u)$ を適用\\[4pt]
  $F[e^{-a^{2}x^{2}} * e^{-b^{2}x^{2}}] = (\pi/(ab))\cdot e^{-u^{2}(1/(4a^{2})+1/(4b^{2}))} = (\pi/(ab))e^{-u^{2}(a^{2}+b^{2})/(4a^{2}b^{2})}$\\[4pt]
  $(2)$ 結果を $\sqrt{\pi/c}\cdot e^{-u^{2}/(4c)}$ の形と比較 $: c = a^{2}b^{2}/(a^{2}+b^{2})$\\[4pt]
  係数を合わせて $e^{-a^{2}x^{2}}*e^{-b^{2}x^{2}} = (\sqrt{\pi}/\sqrt{a^{2}+b^{2}})\cdot e^{-(a^{2}b^{2}/(a^{2}+b^{2}))x^{2}}$
  \medskip\textbf{答}：証明終
\end{frame}

\begin{frame}{演習問題}
  \begin{exampleblock}{自分で解いてみよう}
  \begin{description}
    \item[問5] 性質 III および VI を証明せよ
    \item[問6] $F[f(x)] = F(u), F[g(x)] = G(u)$ のとき、次の等式を証明せよ
    $F[f(x)g(x)] = (1/(2\pi)) F(u) * G(u)$\\
    \item[問7] 96ページの性質を用いて、次の関数のフーリエ変換を求めよ
    $(1) xe^{-x^{2}/2}$\\
    $(2) x^{2}e^{-x^{2}/2}$\\
    \item[問8] 次の等式を証明せよ
    $F^{-1}[e^{-bu^{2}}] = 1/(2\sqrt{\pi b}) \cdot e^{-x^{2}/(4b)}$（b は正の定数）\\
    \item[問9] 次の等式を証明せよ
    $(1) F[e^{-x^{2}/2} * xe^{-x^{2}/2}] = -2\pi iu\cdot e^{-u^{2}}$\\
    $(2) e^{-x^{2}/2} * xe^{-x^{2}/2} = (\sqrt{\pi}/2)\cdot xe^{-x^{2}/4}$\\
  \end{description}
  \end{exampleblock}
\end{frame}

\end{document}
